% Generated draft: requires mathematical and source review.
\documentclass[11pt]{article}
\usepackage[margin=25mm]{geometry}
\usepackage{fontspec}
\setmainfont{Latin Modern Roman}
\usepackage{amsmath,amssymb,mathtools,unicode-math}
\setmathfont{Latin Modern Math}
\usepackage{enumitem,needspace,xurl,hyperref}
\setlength{\parindent}{0pt}
\setlength{\parskip}{5pt}
\setlength{\emergencystretch}{3em}
\newcommand{\N}{\mathbb{N}}
\newcommand{\Z}{\mathbb{Z}}
\newcommand{\Q}{\mathbb{Q}}
\newcommand{\R}{\mathbb{R}}
\newcommand{\C}{\mathbb{C}}
\newcommand{\F}{\mathbb{F}}
\newcommand{\A}{\mathbb{A}}
\newcommand{\PP}{\mathbb{P}}
\newcommand{\E}{\mathbb{E}}
\newcommand{\eps}{\varepsilon}
\newcommand{\dd}{\,\mathrm{d}}
\begin{document}
\section*{UnsolvedMath ID: 800997}
\textbf{Source problem code:} catalogue entry
\subsection{Definitions and mathematical statement}
The supplied record asks for a classification of quantifier patterns at the fourth level of the arithmetical hierarchy. Write \(\mathcal P_{\Sigma^0_4}\) and \(\mathcal P_{\Pi^0_4}\) for the source-intended collections of \(\Sigma^0_4\) and \(\Pi^0_4\) quantifier patterns, respectively. These symbols are placeholders: the supplied record does not define what constitutes a quantifier pattern or specify the relations denoted by \(\equiv_m\) and \(\equiv_{dm}\).

Provided the intended collections and equivalence relations are specified, write
\[
[P]_{\sim}=\{Q\in\mathcal P:Q\sim P\},\qquad
\mathcal P/{\sim}=\{[P]_{\sim}:P\in\mathcal P\}.
\]
The requested problems are:
\begin{enumerate}
\item Determine the number of \(\equiv_m\)-equivalence classes of \(\Sigma^0_4\) quantifier patterns, namely
\[
\left|\mathcal P_{\Sigma^0_4}/{\equiv_m}\right|.
\]
\item Determine the number of \(\equiv_m\)-equivalence classes of \(\Pi^0_4\) quantifier patterns, namely
\[
\left|\mathcal P_{\Pi^0_4}/{\equiv_m}\right|.
\]
\item Determine the number of \(\equiv_{dm}\)-equivalence classes of \(\Sigma^0_4\) quantifier patterns, namely
\[
\left|\mathcal P_{\Sigma^0_4}/{\equiv_{dm}}\right|.
\]
\end{enumerate}
No finiteness assumption is supplied. The meanings of the pattern collections and both equivalence relations must be recovered from the source before these questions have a fully specified mathematical interpretation. In particular, no interpretation of the subscripts \(m\) or \(dm\) is imposed here.
\subsection{Short English statement}
Determine how many equivalence classes the fourth-level arithmetical quantifier patterns have: for \(\Sigma^0_4\) and \(\Pi^0_4\) under \(\equiv_m\), and for \(\Sigma^0_4\) under \(\equiv_{dm}\). The supplied excerpt does not define the patterns or the equivalence relations, so this draft retains those unresolved source-dependent terms.
\subsection{Sources}
\begin{itemize}
\item Local UnsolvedMath registry, permanent ID \(800997\), exact source problem code \texttt{catalogue entry}, titled \emph{Open Problems in Mathematical Logic --- Problem 9}. The supplied registry labels the problem \emph{open}; this is a supplied status, not an independently verified mathematical assessment.
\item Supplied source attribution: Barmpalias, Gao et al., \emph{Open Problems in Mathematical Logic} (2026), Problem 9, \url{https://arxiv.org/abs/2608.26628}. The registry also lists Barmpalias, Gao et al. in its proposer field, with proposed year 2026; the excerpt does not distinguish collection authorship from original proposal attribution.
\item Supplied provenance: source accessed 2026-10-16; extraction method \texttt{arxiv\_latex}; upstream revision \texttt{d371ee917ac9647fbb06cd6fc529deec6f3dda82}. The access date and provenance are supplied metadata; no independent source visit or verification is asserted.
\end{itemize}
\end{document}
